\documentclass[11pt]{article}
\usepackage[a4paper,margin=27mm]{geometry}
\usepackage{amsmath,amssymb,amsthm,booktabs,array,longtable}
\usepackage[T1]{fontenc}
\usepackage{lmodern}
\usepackage[hidelinks]{hyperref}
\newtheorem{theorem}{Theorem}
\newtheorem{lemma}[theorem]{Lemma}
\newtheorem{corollary}[theorem]{Corollary}
\newcommand{\fl}[1]{\left\lfloor #1\right\rfloor}
\newcommand{\J}{\mathcal J}
\newcommand{\ph}{\varphi}
\title{Refutation and correction of Theorem 2.10 (Conjectured),\texorpdfstring{\\}{ }arXiv:1404.0484v1, Section 2, page 7:\texorpdfstring{\\}{ }extensions of Patel (2026) and Raso--Kok (2026)}
\author{Haobo Ma\\\small ChronoAI Pte Ltd, Singapore
\and Wenlin Zhang\\\small National University of Singapore\\\small ChronoAI Pte Ltd, Singapore}
\date{4 November 2026}
\hypersetup{pdftitle={Refutation and correction of Theorem 2.10 (Conjectured), arXiv:1404.0484v1, Section 2, page 7: extensions of Patel (2026) and Raso--Kok (2026)},pdfauthor={Haobo Ma and Wenlin Zhang}}
\begin{document}
\maketitle
\begin{abstract}
Extending the degree formulas of Patel (August 2026) and Raso--Kok
(1 November 2026), we refute and correct Theorem 2.10 (Conjectured) of
Kok, Fisher, Wilkens, Mabula and Mukungunugwa, arXiv:1404.0484v1,
Section 2, page 7: the proposed addition
law for maximum degrees of finite Jaco graphs fails at the distinct orders
$3$ and $11$, even though both graphs have a unique maximum-degree vertex.
Its other branch fails at orders $4$ and $10$. A floor formula gives the
exact replacement: the addition defect is $-1$, $0$ or $1$, determined by
two fractional parts rather than uniqueness. We also prove the same
source's conjectured Lemma 2.2, give the partition requested on page 11,
and correct the exceptional Fibonacci sums in Proposition 3.2.
Patel's maximum-set identities supply the missing order condition in
Theorem 2.3 and correct the edge assertion of Theorem 2.15.
We give the recurrence form of Kok's 2016 Conjecture 2.16.1 and explain
that its McPherson nonrepetition Conjecture 2.48.1 follows directly from
the same article's Theorem 2.48 and the endpoint formula.
\end{abstract}
\begin{center}\small Licensed under Creative Commons Attribution 4.0 International
(CC BY 4.0; \texttt{cc-by-4.0}).\end{center}

\section{The conjecture and the recursive graph}
Write $\Delta_n=\Delta(J_n(1))$ for the ordinary maximum degree of the
underlying simple graph, and $\J_n$ for its set of maximum-degree vertices.
The original statement \cite[Theorem 2.10, p.~7]{original}, labelled
``Conjectured'', asserts, for $n,m\ge3$ and $n\ne m$, that
\begin{equation}\label{posed}
 \Delta_{n+m}=\begin{cases}
 \Delta_n+\Delta_m,& |\J_n|=1\text{ or }|\J_m|=1,\\
 \Delta_n+\Delta_m+1,&\text{otherwise}.
 \end{cases}
\end{equation}
The same paper repeats this as an open problem on page 11.

We use precisely its Definitions 1.1--1.3. The vertices have positive
labels $1,\ldots,n$. Let $a_i$ be the infinite indegree and $R_i$ the last
outgoing label. Recursively,
\begin{equation}\label{recurrence}
 a_i=\#\{j:1\le j<i,\ R_j\ge i\},\qquad R_i=2i-a_i.
\end{equation}
For $i<j$, the arc $v_i\to v_j$ exists exactly when $j\le R_i$.
Degrees below count both incoming and outgoing edges, not outdegrees alone.
Put $\ph=(1+\sqrt5)/2$ and $b_i=\fl{(i+1)/\ph}$ for $i\ge1$.

\begin{lemma}\label{formula}
For all positive labels and all $1\le i\le n$,
\begin{equation}\label{degree}
 R_i=i+b_i=\fl{(i+1)\ph}-1,\qquad
 \deg_{J_n(1)}(v_i)=\min\{i,n-b_i\}.
\end{equation}
Consequently, for every $n\ge1$,
\begin{equation}\label{maximum}
 \Delta_n=\fl{\frac{n+2}{\ph}}-1.
\end{equation}
\end{lemma}
\begin{proof}
The candidate $R_j^*=\fl{(j+1)\ph}-1=j+b_j$ satisfies
\[
 i\le R_j^*\quad\Longleftrightarrow\quad
 i+1\le\fl{(j+1)\ph}\quad\Longleftrightarrow\quad b_i\le j.
\]
The last equivalence uses irrationality: $(i+1)/\ph<j+1$ is equivalent
to $\fl{(i+1)/\ph}\le j$. Since $1\le b_i\le i$, strong induction in
\eqref{recurrence} identifies the incoming labels as $b_i,\ldots,i-1$,
so $a_i=i-b_i$ and $R_i=i+b_i$. There are $\min(b_i,n-i)$ outgoing
neighbors. This proves the degree formula.

Monotonicity of $b_i$ gives, for $1\le r\le n$,
\[
 \Delta_n\ge r\quad\Longleftrightarrow\quad
 \exists i\ge r:\ b_i\le n-r
 \quad\Longleftrightarrow\quad r+b_r\le n.
\]
This is $\fl{(r+1)\ph}\le n+1$, equivalently $r+1<(n+2)/\ph$.
For $n\ge2$ the largest such integer is \eqref{maximum}; it lies in
$[1,n]$ since $4/\ph>2$ and $(n+2)/\ph<n+2$. For $n=1$ the graph
has degree zero, also given by \eqref{maximum}.
\end{proof}

Patel's MathDB solution \cite{patel} proves these finite degree and
maximum formulas. Raso and Kok \cite[Theorems 4.5 and 4.8]{raso}
give a later independent proof, posted on 1 November 2026; their
reference 7 dates Patel's solution to August 2026. We include the interval
argument to make the addition rule and corrections self-contained.
The same floor expression also follows by combining the original
Fisher Algorithm, Steps 2--3, with Bettina's Theorem 3.1 in
\cite{original} and the classical Hofstadter identity recorded in
Rahman \cite[Section 1, p.~1]{rahman}: the Zeckendorf down-shift is
$G(N)=\fl{(N+1)/\ph}$, so $\Delta_n=G(n+1)-1$ for $n\ge3$.
This identifies earlier ingredients for the formula; the addition
defect below is its exact shifted-floor consequence.
Kok's journal account \cite[Conjectures 2.4.1--2.4.2 and
Proposition 2.21, pp.~72--74]{critical} repeats the conjectures and
Fibonacci assertions. Its Table 1 already contains the numerical values
used below; the issue is their compatibility with the stated laws.
Its Section 2.7, page 83, records Wagner's earlier indegree formula.

\section{Counterexamples and the exact addition rule}
Equations \eqref{degree}--\eqref{maximum}, or the integer recursion
\eqref{recurrence} directly, give the following values. Labels in the last
column denote vertices.
\begin{center}
\begin{tabular}{rcl}
\toprule
$n$ & $\Delta_n$ & $\J_n$\\
\midrule
3 & 2 & $\{2\}$\\
4 & 2 & $\{2,3\}$\\
6 & 3 & $\{3,4,5\}$\\
10 & 6 & $\{6,7\}$\\
11 & 7 & $\{7\}$\\
14 & 8 & $\{8,9,10\}$\\
\bottomrule
\end{tabular}
\end{center}
In particular $3,11\ge3$ and $3\ne11$, and both $\J_3$ and $\J_{11}$
are singletons, but
\[
 \Delta_{14}=8\ne2+7=\Delta_3+\Delta_{11}.
\]
Thus even requiring uniqueness in \emph{both} summands does not repair
the first branch. Neither $\J_4$ nor $\J_{10}$ is a singleton, whereas
$\Delta_{14}=8\ne2+6+1=9$, disproving the second branch too.

\begin{theorem}\label{addition}
Let $n,m\ge1$, set $x_n=\{(n+2)/\ph\}$ and $\alpha=2/\ph-1=\sqrt5-2$.
Then
\begin{equation}\label{corrected}
 \Delta_{n+m}-\Delta_n-\Delta_m
 =\fl{x_n+x_m-\alpha}
 =\begin{cases}
 -1,&x_n+x_m<\alpha,\\
 0,&\alpha\le x_n+x_m<1+\alpha,\\
 1,&1+\alpha\le x_n+x_m.
 \end{cases}
\end{equation}
All three values occur for distinct orders at least three.
\end{theorem}
\begin{proof}
Put $u=(n+2)/\ph$ and $v=(m+2)/\ph$. Formula \eqref{maximum} gives
\[
 \Delta_{n+m}-\Delta_n-\Delta_m
 =\fl{u+v-2/\ph}-\fl u-\fl v+1
 =\fl{\{u\}+\{v\}-\alpha}.
\]
Since $0\le x_n+x_m<2$ and $0<\alpha<1$, the three cases follow.
The pairs $(3,11)$, $(3,4)$ and $(4,6)$ give respectively $-1,0,1$.
\end{proof}

\begin{corollary}\label{many}
For $r\ge2$ and positive orders $n_1,\ldots,n_r$, put
$N=\sum_{j=1}^r n_j$ and $x_j=\{(n_j+2)/\ph\}$. Then
\[
 \Delta_N-\sum_{j=1}^r\Delta_{n_j}
   =\fl{\sum_{j=1}^r x_j-(r-1)\alpha}.
\]
\end{corollary}
\begin{proof}
Sum $(n_j+2)/\ph$ and subtract $2(r-1)/\ph$ to obtain
$(N+2)/\ph$. Apply \eqref{maximum} and move the integer $r-1$
inside the remaining floor, using $2/\ph=1+\alpha$.
\end{proof}

\section{Adjacent source claims using known maximum-set identities}
We use Patel's maximum-set description \cite{patelsets} as a cited tool,
not as an additional classification result. The related maximum-set
questions are Kok's Conjectures 2.3--2.6 in \cite[Section 2.2]{sequences}.
Here we retain the identities needed for the neighboring 2014 claims.
For $n\ge2$, put $p=\Delta_n$ and $t=n-p$. Patel's identities give
\begin{equation}\label{setinterval}
 \J_n=\{v_i:p\le i\le\min(n,R_t)\},\qquad 1\le|\J_n|\le3,
\end{equation}
and
\begin{equation}\label{uniquecriterion}
 \J_n=\{v_p\}\quad\Longleftrightarrow\quad
 n=R_p\ \hbox{and}\ b_{p+1}=b_p+1.
\end{equation}
For clarity, uniqueness is also equivalent to $b_{n+1}=b_{n+2}$:
if the maximum is not unique, $v_{p+1}$ has degree $p+1$ at order
$n+1$, whereas if it is unique every label above $p$ has degree at
most $p$ there. Labels at most $p$ always have degree at most $p$,
and larger labels gain at most one edge. Thus
$\Delta_{n+1}=\Delta_n$ exactly in the unique case. Formula
\eqref{maximum} makes this $b_{n+2}=b_{n+1}$.

\paragraph{Consequences for the source's Section 2.}
The bound proves Proposition 2.4. Criterion \eqref{uniquecriterion}
proves Proposition 2.1 and Corollary 2.6: at a unique maximum,
$R_p=n$ and $R_{p+1}=n+2$. It also proves Proposition 2.7, since
equal indegrees at $k$ and $k+1$ mean $b_{k+1}=b_k+1$.
The last equivalence proves Proposition 2.8 for $k\ge3$ after
putting $n=k-1$; its $k=2$ boundary is false because $\J_1=\{v_1\}$
but $b_2\ne b_3$. At a unique maximum $p$ the next maximum stays
$p$ and has two or three labels, proving Corollary 2.5 for $n\ge2$.
Conversely, if the next maximum set starts at the same label $p$,
the maximum did not increase, so the preceding maximum was unique.
The singleton graph requires separate treatment.
Indeed $\J_1=\{v_1\}$ and $\J_2=\{v_1,v_2\}$ directly.

Theorem 2.11 follows for $n\ge2$: uniqueness gives $R_p=n$,
so $v_pv_n$ is an edge and
$\deg(v_n)=n-b_n=n-p$. Indeed \eqref{uniquecriterion} and
the endpoint inversion give $b_n=p$.
More generally, for $n\ge1$, Theorem 2.9 has the refinement
\begin{equation}\label{nested}
 \Delta_{n+\Delta_n}=
 \begin{cases}n,&\Delta_n=b_n,\\ n-1,&\Delta_n=b_n-1.\end{cases}
\end{equation}
These are the only possibilities because
$\Delta_n=b_{n+1}-1$. The new order is respectively $R_n$ or
$R_n-1$. The threshold then gives respectively maximum $n$ or $n-1$;
the case $n=1$ is direct.

\begin{corollary}\label{morrie}
Theorem 2.3 of \cite{original} becomes valid for $n\ge2$ with prime
maximum label $k$ after adding $n=R_k$ to each left-hand condition.
Its unqualified parts (a) and (b) fail respectively at $(n,k)=(6,3)$
and $(4,2)$.
\end{corollary}
\begin{proof}
Write $a_i=i-b_i$. Under $n=R_k$ and $a_{k+1}=a_k$,
criterion \eqref{uniquecriterion} gives $\J_n=\{v_k\}$.
At $n+1$, labels $k,k+1$ are maximum. Label $k+2$ is also maximum
exactly when $b_{k+2}=b_k+1$, equivalently $a_{k+2}=a_k+1$;
otherwise $b_{k+2}=b_k+2$ and $a_{k+2}=a_k$. The cardinality
bound excludes further labels. Conversely, the asserted singleton
at $n$ forces $n=R_k$ and $a_{k+1}=a_k$, and the next maximum
set determines the indicated value of $a_{k+2}$.
For the counterexamples, $a_3=a_4=1$, $a_5=2$, but
$\J_6=\{v_3,v_4,v_5\}$; also $a_2=a_3=a_4=1$, but
$\J_4=\{v_2,v_3\}$.
\end{proof}

\begin{corollary}\label{edgecorrection}
For $n=7+3k$, $k\ge0$, the bound in Theorem 2.15 holds:
$\Delta_n\le4+2k=n-(3+k)$. Its claimed edge $v_{\Delta_n}v_n$
exists exactly when $n=R_{\Delta_n}$, and fails at $n=19$.
\end{corollary}
\begin{proof}
Since $9+3k<(6+2k)\ph$, formula \eqref{maximum} gives the bound.
The threshold gives $R_{\Delta_n}\le n$, while the edge requires
the reverse inequality. At $n=19$ we have $\Delta_{19}=11$ and
$R_{11}=18$, so the edge is absent.
\end{proof}

\section{The neighboring conjecture and open-problem list}
\begin{corollary}\label{tripleorders}
The recurrence form of Conjecture 2.16.1 of \cite[p.~73]{critical}
is as follows. Set $p_0=1$ and, for $j\ge0$, define
\[
 p_{j+1}=p_j+3+2\left(\fl{(j+2)/\ph}-\fl{(j+1)/\ph}\right).
\]
The actual orders with three maximum-degree vertices are
$p_1,p_2,\ldots=6,9,14,19,22,\ldots$.
Successive increments are $3$ or $5$, with no consecutive $3$s
and no three consecutive $5$s. The order-one value is a synthetic
seed, not a graph with three maximum-degree vertices.
This is the recurrence reformulation of Patel's prior classification
\cite{patelsets}; the full maximum-set and never-singleton-label
classifications \cite{patelsets,patelnever} are separate results.
\end{corollary}
\begin{proof}
Subtract consecutive terms in the three-maximizer case of Patel's
classification \cite{patelsets}; using $\ph^2=2+1/\ph$ gives
exactly the displayed recurrence, with first actual order $6$.
Its first gap from the synthetic seed is $5$, as required.
The floor increments are zero or one.
Two zero increments are impossible because $2/\ph>1$;
three unit increments are impossible because $3/\ph<2$.
\end{proof}

\begin{corollary}\label{mcpherson}
Conjecture 2.48.1 of \cite[p.~91]{critical} follows directly from
its Theorem 2.48 and the endpoint formula: for $n\ge3$,
the infinite outdegree $b_n$ is nonrepetitive exactly when the
McPherson number $\Upsilon_n$ is nonrepetitive. More precisely,
$\Upsilon_n=b_n-1$ for every $n\ge1$.
\end{corollary}
\begin{proof}
For $n\ge3$, write $p=\Delta_n$, the prime maximum label.
Theorem 2.48 of \cite{critical} already gives $\Upsilon_n=p-1$
if $v_pv_n$ is an edge, and $\Upsilon_n=p$ otherwise.
The same formula appears earlier in Kok and Susanth
\cite[Theorem 2.3, p.~4]{mcphersonprior}, posted 31 October 2014;
their page 5 already poses the nonrepetition conjecture.
The maximum threshold is $R_p\le n<R_{p+1}$, and endpoint inversion
gives $b_n=p$ in the edge case $n=R_p$, and $b_n=p+1$ in the
other case. Indeed $R_{p-1}<n$ and $R_{p+1}>n$, so $b_n$ is
either $p$ or $p+1$, with the former exactly when $R_p\ge n$.
Thus the two branches of the cited theorem are precisely
$\Upsilon_n=b_n-1$. Orders $1$ and $2$ satisfy this directly.
Subtracting one preserves both strict inequalities defining
nonrepetition, yielding the stated equivalence.
\end{proof}

\begin{corollary}\label{nonrepeat}
Conjectured Lemma 2.2 of \cite[p.~4]{original} holds on its positive-label
domain: if $n\ge2$ and $b_{n-1}<b_n<b_{n+1}$, then
$\J_n=\{v_{b_n}\}$.
\end{corollary}
\begin{proof}
Consecutive $b_i$ differ by zero or one. Write $\ell=b_n$; the hypotheses
give $b_{n-1}=\ell-1$ and $b_{n+1}=\ell+1$. Thus
$\Delta_n=b_{n+1}-1=\ell$. Moreover $\fl{n/\ph}=\ell-1$ implies
$n<\ell\ph$, whence
\[
 \frac{\ell+2}{\ph}>n-\ell+1\qquad(2/\ph>1).
\]
Therefore $b_{\ell+1}\ge n-\ell+1$ and every label $i\ge\ell+1$
has degree at most $n-b_i\le\ell-1$. Every label $i<\ell$ has degree
at most $i<\ell$. The maximum $\ell$ is therefore attained only at $v_\ell$.
\end{proof}

The converse, which is not asserted in the source, fails: at $n=5$ the maximum
is unique, but $b_4=b_5=3$.

\begin{corollary}\label{partition}
The first open problem on page 11 of \cite{original} has the following
partition. Set $\beta=5/\ph-3>0$, $k_0=1$ and, for $q\ge1$, set
\begin{equation}\label{boundaries}
 k_q=\left\lceil\frac{q-\alpha}{\beta}\right\rceil
     =\fl{(5q-1)\ph}+3q.
\end{equation}
Take $P_0=\{1,\ldots,5k_1-1\}$ and
$P_q=\{5k_q,\ldots,5k_{q+1}-1\}$ for $q\ge1$. These partition the
positive integers, and whenever $n=5k\in P_q$,
\begin{equation}\label{five}
 \Delta_{5k}=3k+q.
\end{equation}
In particular $P_0=[1,44]$, $P_1=[45,99]$ and $P_2=[100,154]$.
\end{corollary}
\begin{proof}
Equation \eqref{maximum} gives
$\Delta_{5k}-3k=\fl{\beta k+\alpha}$.
This equals $q$ precisely when $k_q\le k<k_{q+1}$, with the lower
bound restricted to $k\ge1$ at $q=0$. Rationalizing the quotient yields
$(q-\alpha)/\beta=(5q-1)\ph+3q-1$; irrationality gives
\eqref{boundaries}. Since $0<\beta<1$, the $k_q$ strictly increase
and tend to infinity, proving the partition assertion.
\end{proof}

The same argument gives every arithmetic progression: for $d\ge1$,
put $c=\fl{d/\ph}$ and $\beta_d=d/\ph-c\in(0,1)$. Then
\[
 \Delta_{dk}=ck+\fl{\beta_d k+\alpha}\qquad(k\ge1).
\]
The part with correction $q$ is exactly
$\max(1,\lceil(q-\alpha)/\beta_d\rceil)\le k
<\lceil(q+1-\alpha)/\beta_d\rceil$ for $q\ge0$.
This follows by substituting $n=dk$ in \eqref{maximum}.

\section{The Fibonacci claims that use the conjecture}
Use $F_0=0$, $F_1=1$, $F_{r+1}=F_r+F_{r-1}$. Proposition 3.2 of
\cite[pp.~10--11]{original} takes consecutive Fibonacci numbers
$t=F_{r-1}$, $h=F_r$ and another Fibonacci number $l=F_s$, with
$t,h,l\ge3$. Thus its literal domain is $r\ge5$, $s\ge4$.
It states that $\Delta_h=F_{r-1}$ and $\J_h=\{v_{F_{r-1}}\}$,
and that $\Delta_{h+l}=\Delta_h+\Delta_l$ and
$\J_{h+l}=\{v_{\Delta_h+\Delta_l}\}$. Its proofs of the last two
assertions invoke the conjectured Theorem 2.10. The next result also
gives the symmetric extension $r,s\ge4$, distinguishing its extra boundary.

\begin{corollary}\label{fib}
The first three assertions of Proposition 3.2 hold throughout its stated
domain. Its singleton-sum assertion fails exactly at $(h,l)=(5,5)$.
More generally, for $r,s\ge4$, the sum assertions are
\[
 \Delta_{F_r+F_s}=F_{r-1}+F_{s-1},\qquad
 \J_{F_r+F_s}=\{v_{F_{r-1}+F_{s-1}}\},
\]
except that the degree assertion fails at $(F_r,F_s)=(3,3)$, and the
singleton assertion fails at $(3,3)$ and $(5,5)$. The pair $(3,3)$ is
outside the literal domain of Proposition 3.2 because its $t$ would be two.
At these exceptions, $\Delta_6=3$, $\J_6=\{v_3,v_4,v_5\}$ and
$\Delta_{10}=6$, $\J_{10}=\{v_6,v_7\}$. Thus both assertions hold
without exceptions if the two Fibonacci orders are distinct.
\end{corollary}
\begin{proof}
The Fibonacci recurrence gives
\[
 F_r/\ph=F_{r-1}+\varepsilon_r,\qquad
 \varepsilon_r=(-1)^{r-1}\ph^{-r}.
\]
For $r\ge4$, $-\ph^{-4}\le\varepsilon_r\le\ph^{-5}$; hence
$1<2/\ph+\varepsilon_r<2$. Formula \eqref{maximum} proves
$\Delta_{F_r}=F_{r-1}$.

For any $n\ge2$ with $\Delta_n=p$, equation \eqref{degree} shows
that $\J_n=\{v_p\}$ exactly when $b_{p+1}>n-p$:
labels below $p$ have smaller degree, while labels above $p$ do so
exactly when $v_{p+1}$ does. Writing $n/\ph=p+E$, this condition is
equivalent to
\[
 \frac{p+2}{\ph}>n-p+1
 \quad\Longleftrightarrow\quad E<\frac{2-\ph}{\ph^2}=\ph^{-4}.
\]
For $n=F_r$, $E=\varepsilon_r<\ph^{-4}$, proving the singleton assertion.

For a sum put $p=F_{r-1}+F_{s-1}$ and $E=\varepsilon_r+\varepsilon_s$.
Apart from $r=s=4$,
\[
 -\ph^{-4}-\ph^{-6}\le E\le2\ph^{-5},\qquad
 \ph^{-4}+\ph^{-6}<2/\ph-1,\quad 2\ph^{-5}<2-2/\ph.
\]
Consequently $1<2/\ph+E<2$ and $\Delta_{F_r+F_s}=p$.
For $r=s=5$, $E=2\ph^{-5}>\ph^{-4}$, so the maximum is not unique.
For every remaining pair other than $r=s=4$, either both errors are
positive with at most one equal to $\ph^{-5}$, or at least one is
negative. Thus $E\le\ph^{-5}+\ph^{-7}<\ph^{-4}$, giving uniqueness.
All the displayed comparisons follow from $\ph^2=\ph+1$ and
$8/5<\ph<13/8$. The two exceptional graphs are in the preceding table.
\end{proof}

\section{Coverage and verification}
The source-item coverage is as follows; the locators refer to
arXiv:1404.0484v1, not to a renumbered later article.
\begingroup\small
\begin{longtable}{@{}>{\raggedright\arraybackslash}p{.35\linewidth}>{\raggedright\arraybackslash}p{.60\linewidth}@{}}
\toprule
Source item & Mathematical coverage\\
\midrule
\endhead
Theorem 2.10, p.~7; third open problem, p.~11 & Both branches refuted;
refuted-and-corrected by Theorem~\ref{addition};
Corollary~\ref{many} extends it to any number of summands.\\
Degree and maximum formulas & Reproved; prior proof Patel (2026),
later independent proof Raso--Kok (2026) \cite{patel,raso}.\\
Maximum-set intervals & Prior proof cited: Patel (2026)
\cite{patelsets}; used as tools in \eqref{setinterval}--\eqref{uniquecriterion}.
The separate maximum-set questions are \cite[Conjectures 2.3--2.6]{sequences}.\\
Lemma 2.2, p.~4; second open problem, p.~11 & Proved by
Corollary~\ref{nonrepeat}; converse (not asserted in the source) fails at $n=5$.\\
First open problem, p.~11 & Partition given in
Corollary~\ref{partition}, including endpoint $154$; extended to
every arithmetic progression.\\
Proposition 2.1; Corollary 2.6; Proposition 2.7 & Reproved; prior
proof Kok et al. (2014). Criterion \eqref{uniquecriterion} gives
the assertions on their positive-label domains.\\
Theorem 2.3(a)--(b), pp.~4--5 & Refuted-and-corrected:
the missing condition is $n=R_k$; Corollary~\ref{morrie}.\\
Proposition 2.4; Corollary 2.5 & Reproved; prior proof Kok et al.
(2014). Cardinalities and transitions follow from
\eqref{setinterval}--\eqref{uniquecriterion}; singleton graph treated separately.\\
Proposition 2.8 & Refuted-and-corrected boundary: valid for $k\ge3$,
false at $k=2$; \eqref{setinterval}--\eqref{uniquecriterion}.\\
Theorems 2.9 and 2.11 & Reproved; prior proof Kok et al. (2014).
Equation \eqref{nested} selects the exact alternative in 2.9;
2.11 holds for $n\ge2$.\\
Theorem 2.15, p.~9 & Refuted-and-corrected: bound retained, edge
criterion supplied; edge absent at $n=19$, Corollary~\ref{edgecorrection}.\\
\cite{critical}, Conjectures 2.4.1--2.4.2; Proposition 2.21 &
Same resolutions as the corresponding 2014 items above; later
posing locators, not earlier resolutions.\\
\cite{critical}, Conjecture 2.48.1, p.~91 & Prior theorem cited:
follows directly from Theorem 2.48 of Kok (2016) and the endpoint formula;
the formula is already in Kok--Susanth (2014), Theorem 2.3
\cite{mcphersonprior}. Corollary~\ref{mcpherson} spells out the implication.\\
\cite{critical}, Conjecture 2.16.1, p.~73 & Recurrence form only;
prior classification Patel (2026) \cite{patelsets};
exact recurrence and seed boundary in
Corollary~\ref{tripleorders}.\\
Proposition 3.2(a)--(d), pp.~10--11 & (a)--(c) proved; (d) corrected
at $5+5$. Corollary~\ref{fib} also covers the extended boundary $3+3$.\\
Theorem 3.1 (Bettina), pp.~9--10 & Prior proof cited in
\cite{original}; its Fibonacci outdegree representation is distinct
from the maximum-degree addition rule.\\
Spanning-tree enumeration, fourth open problem, p.~11 & Not covered.\\
Other introductory, algorithmic and adjacency statements & Not
separately covered; the recursive definition is used directly.\\
Other parameters and general slopes in \cite{critical} & Not covered.\\
\bottomrule
\end{longtable}
\endgroup

The pinned Lean declaration
\href{https://github.com/the-omega-institute/trureturing/blob/3bedf2870aacc6e7f17353444aadaa49fb73157c/D5/S3/ConceptDynamics/GraphColoring/JacoAdditionLawRefutation.lean}{\texttt{JacoAdditionLawRefutation.result}}
states the negation of the full two-branch assertion \eqref{posed}; its
proof uses orders $3$ and $11$ and establishes uniqueness at order $3$.
Uniqueness at order $11$ has the written and finite proofs given here.
In that file, \texttt{UniqueJaconian}
means exactly one vertex attaining the ordinary maximum degree, and
\texttt{claim} includes $n,m\ge3$ and $n\ne m$.
The same-revision file
\href{https://github.com/the-omega-institute/trureturing/blob/3bedf2870aacc6e7f17353444aadaa49fb73157c/D5/S3/ConceptDynamics/GraphColoring/JacoMaximumDegree.lean}{\texttt{JacoMaximumDegree.lean}}
contains \texttt{degree\_formula} and \texttt{result}, which give
\eqref{degree} and \eqref{maximum} (the floor clause of \texttt{result}
has $n\ge2$). These declarations use the induced graph on labels $1,\ldots,n$.
The same-revision
\href{https://github.com/the-omega-institute/trureturing/blob/3bedf2870aacc6e7f17353444aadaa49fb73157c/D5/S3/ConceptDynamics/GraphColoring/JacoEndpointBeatty.lean}{\texttt{JacoEndpointBeatty.lean}}
contains \texttt{endpoint\_recurrence} and \texttt{endpoint\_eq\_beatty},
connecting recursive counting to the endpoint $i+b_i$.
At that same revision, the linked
\href{https://github.com/the-omega-institute/trureturing/blob/3bedf2870aacc6e7f17353444aadaa49fb73157c/D5/S3/ConceptDynamics/GraphColoring/JacoDomPathDiameterRefutation.lean}{\texttt{Vertex} and \texttt{jaco} definitions}
select labels $1,\ldots,n$ and their induced graph.
The \href{https://github.com/the-omega-institute/trureturing/blob/3bedf2870aacc6e7f17353444aadaa49fb73157c/D5/S0/Certificates/JacoExponentialDominationRefutation.lean}{\texttt{Adj} definition}
uses the recursive endpoint array.
The second-branch counterexample and Theorem~\ref{addition} have written proofs.
The corollaries have written arguments, using the explicitly cited prior
maximum-set classification and McPherson theorem where indicated.
The accompanying \texttt{verify.py} constructs the recursive graphs and
checks every degree for $1\le i\le n\le512$, all distinct pairs
$3\le n<m\le256$, the nonrepetitive-outdegree implication through
$n=511$, the partition through $5k\le512$, and all Fibonacci pairs
of orders at least three whose sum is at most $512$.
It also checks maximum-set intervals and the neighboring claims through
order $512$, the three-maximizer recurrence through that order,
three-summand additions with summands at most $32$,
and every progression $dk\le512$ with $1\le d\le32$.
Clique intervals are checked through order $512$, and all exploded
subsets and cliques are enumerated for orders at most $12$.

\begin{thebibliography}{9}
\bibitem{original} J.~Kok, P.~Fisher, B.~Wilkens, M.~Mabula and
V.~Mukungunugwa, \emph{Characteristics of Finite Jaco Graphs,
$J_n(1)$, $n\in\mathbb N$},
\href{https://arxiv.org/abs/1404.0484v1}{arXiv:1404.0484v1}, 2 April 2014.
\bibitem{patel} Shivam Patel, solution to \emph{Maximum-degree formula
conjecture for finite linear Jaco graphs},
\href{https://mathdb.com/p/369423/maximum-degree-formula-conjecture-for-finite-linear-jaco-gra}{MathDB problem 369423},
2026 (August date reported in \cite[reference 7]{raso}).
\bibitem{patelsets} Shivam Patel, solution to \emph{Delta-set cardinality
conjecture for finite linear Jaco graphs},
\href{https://mathdb.com/p/369424/delta-set-cardinality-conjecture-for-finite-linear-jaco-grap}{MathDB problem 369424},
2026.
\bibitem{patelnever} Shivam Patel, solution to \emph{Vertices that never
form singleton Delta-sets in finite linear Jaco graphs},
\href{https://mathdb.com/p/369425/vertices-that-never-form-singleton-delta-sets-in-finite-line}{MathDB problem 369425},
2026.
\bibitem{critical} J.~Kok, \emph{Linear Jaco Graphs: A Critical Review},
Journal of Informatics and Mathematical Sciences \textbf{8}(2) (2016),
67--103, published 24 May 2016,
\href{https://doi.org/10.26713/jims.v8i2.402}{doi:10.26713/jims.v8i2.402}.
\bibitem{sequences} J.~Kok, \emph{Integer sequences with conjectured
relation with certain graph parameters of the family of linear Jaco graphs},
\href{https://arxiv.org/abs/2507.16500v1}{arXiv:2507.16500v1},
22 July 2025, Section 2.2, Conjectures 2.3--2.6.
\bibitem{mcphersonprior} J.~Kok and Susanth C.,
\emph{Introduction to the McPherson number, $\Upsilon(G)$ of a simple connected graph},
\href{https://arxiv.org/abs/1410.8637v1}{arXiv:1410.8637v1},
31 October 2014, Theorem 2.3, p.~4, and conjecture, p.~5.
\bibitem{raso} T.~Raso and J.~Kok,
\emph{Short Paper No.~1: Relation between integer sequence A319433 and the
maximum degrees of finite linear Jaco graphs},
\href{https://arxiv.org/abs/2611.01554v1}{arXiv:2611.01554v1},
1 November 2026, Theorems 4.5 and 4.8 and reference 7.
\bibitem{rahman} M.~Rahman,
\emph{A combinatorial interpretation of Hofstadter's G-sequence},
\href{https://arxiv.org/abs/1105.1718v3}{arXiv:1105.1718v3},
16 June 2013, Section 1, p.~1, recording the classical floor and
Zeckendorf down-shift identities.
\end{thebibliography}
\end{document}
